\chapter[Studying Religion before Judging]{How to Study Religion without Rushing to Judge}
\section{A String of Beads}
Begin with a small object: prayer beads.

Wood, glass, agate, perhaps olive pits, threaded into a loop, usually with a larger mother bead. Pinch one, recite something, move it past, take the next. Completing the circle counts one round of what occupies your mind.

Strangely, this object belongs simultaneously to several worlds with little contact.

On Lhasa's Barkhor Street, elderly pilgrims pass 108 Buddhist beads while reciting Avalokiteshvara's mantra. Outside Istanbul or Cairo teahouses, men wind 33 or 99 beads through their hands, repeatedly praising God. In a European Catholic church, an elderly woman kneels holding a 59-bead rosary, arranged in groups of ten for scriptural meditation. Hindu practitioners also use 108 beads.

They do not know one another. Doctrines differ greatly; answers about cosmic origins and life after death contradict. Yet fingers move almost identically: bead, phrase, another bead.

Are they doing the same thing? This is our real question.

If yes, are centuries of religious dispute and even warfare merely misunderstandings? If no, how explain matching movements, half-closed eyes, fingers tightening in danger?

Religion appears daily in news, history, arguments. Ask what exactly it means, or why two practices belong together, and most people, adults included, struggle. This chapter teaches neither belief nor unbelief, but something harder and more useful: \textbf{studying religion without rushing to judge}.

\section{Two Places of Worship on One Street}
Come into a scene.

Tumen Street, Quanzhou, Fujian: narrow, with scooters, snack stalls, grandmothers carrying groceries, midday southern sun polishing the asphalt.

Tonghuai Guanyue Temple stands beside it. Incense arrives before the entrance, thick enough to taste, mixed with firecracker residue. Joss-paper shops thrive, stacks behind glass. Cross the threshold and light drops. Images crowd ahead: red-faced, long-bearded Guan Yu central, other divine generals beside him. Banners and lanterns hang overhead; divination-stick holders, burners, donation boxes cluster. An elderly woman raises incense over her head, bows three times, plants it, then shakes two crescent wooden divination blocks in her palms and throws them. One lands each way. She relaxes: ``Sheng jiao---a favourable answer. He has agreed.''

Less than two minutes along the street stands another building: Qingjing Mosque.

Quiet as another world. Stone gateway, pointed arch, Arabic inscriptions: founded in Northern Song times, among China's oldest surviving mosques. No divine statue or portrait, because Islamic teaching prohibits images of God and prophets. At prayer time, white-capped men remove shoes and enter a mat-covered hall, facing one direction, standing, bowing, prostrating in orderly silence. No incense, divination sticks, donation box. Visitors' voices automatically lower.

A little over a hundred seconds' walk separates smoke, images, and crowds from plain emptiness and low recitation. If both are religion, the word has a remarkable appetite, swallowing near-opposites.

Year Eight student Xiaolin visits with his teacher. Wrinkling his nose in Guanyue Temple, he calls it superstition: what can two wooden blocks decide? At the mosque he murmurs that it is clean, but without incense or bowing before gods, what is the point of daily bowing towards a wall?

The teacher only asks: ``Before deciding whether their worship is right, can you tell me what each is \textbf{doing}?''

Xiaolin is startled. We shall return to the question.

\section{Why Is Religion So Hard to Define?}
Put his common intuition on the table: religion means believing in gods.

Neat, but immediately leaky.

First, some religions place no god at their centre. Early Buddhism, to which contemporary Theravada in Sri Lanka, Thailand, and Myanmar claims greatest fidelity, teaches Four Truths, Eightfold Path, dependent arising: analysis and practice explaining suffering's birth through attachment and its cessation. Buddha is human, not divine, repeatedly denying the relevance of a creator to liberation. Later Mahayana teems with Buddhas and bodhisattvas and resembles a divine temple, you might say. Precisely: \textbf{one tradition differs enormously within itself}. The god-belief ruler cuts Buddhism in half.

Second, some temples enshrine humans. Guan Yu was a Three Kingdoms general, posthumously promoted through centuries from marquis to king to emperor, becoming Lord Guan. Confucius was a teacher repeatedly rebuffed in life. Is honouring a historical person religion? If so, what about bows in memorial halls? If not, how does the grandmother's unmistakable devotion differ from a churchgoer's?

Third, what about practices without temples: geomancers inspecting sites, fortune-tellers reading birth data, youths forwarding lucky carp, students honouring Confucius before exams? Supernatural belief might include all, yet they lack churches, scriptures, monastics. Include everything and religion bursts; exclude it and where lies the boundary?

A deeper problem reverses the view. Today religion seems selectable and declarable: what faith are you? Many historical communities lacked this concept. For most ancient Chinese, Qingming ancestors, Dragon Boat Festival mugwort, Kitchen God worship, and commissioned Buddhist rites were living, not joining a faith. Japanese shukyo, religion, was coined in the nineteenth century to translate a Western word. Previously people did not separate nearby shrines and temples into two religions: births reported at shrines, weddings in shrines or churches, funerals with Buddhist chanting. A life crossed between them without demanding one label.

Definitional difficulty is therefore not simply insufficient reading. The basket religion was woven later, then filled with varied practices worldwide, fitting some and forcing others. Its shape came from early fillers, chiefly European scholars modelling familiar Christianity. A church-shaped basket does not fit shrines, ancestor rites, qigong, yoga; blaming the contents gets things backwards.

Must we burn the word? No. Change its use rather than discard it. That is the coming tool's task.

\section{Inside and Outside the Temple}
Return to Tumen Street and state each voice fully.

\textbf{The grandmother inside.} Asked whether Guan Yu truly governs her grandson's exam, she might not answer directly. Her mother-in-law worshipped here; her grandson faces entrance exams; she can offer little else and seeks peace; Lord Guan honours loyalty and trust, so why not honour him? None of this proves divine existence. Practice connects older generations, care for grandchildren, ordinary ethics of loyalty. Her world is a network of relationships and habits, not propositions. Xiaolin's superstition sweeps away the network, not merely blocks.

\textbf{The outside researcher.} An anthropologist also carries a notebook in the temple. Neither worshipping nor mocking, she records attendance, gender, spending on incense, responses to favourable and unfavourable throws. What function does ritual serve here? Professional discipline requires \textbf{withholding an answer} about Guan Yu's efficacy, not concealing an existing verdict but recognising that observation, interviews, statistics, and comparison cannot reach it. A ruler cannot measure temperature; neither temperature nor ruler is at fault.

\textbf{The younger family member.} The teenage grandson may inwardly groan at Grandma's repeated ritual yet leave her exam-eve incense burning. He disbelieves but cannot bear to disappoint. Countless people inhabit this ambiguous cooperation without belief, among religious studies' most intriguing positions.

\textbf{Differences within tradition.} Beware treating a religion as one person with one mouth. All argue internally. Buddhists disagree whether incense and images skilfully welcome ordinary people or perpetuate attachment Buddha sought to dissolve. Christians dispute prayer to Mary and saints over centuries, wars, divisions. Muslims interpret one scripture differently enough to form schools. Hearing any religion thinks something should trigger which period, region, school, person?

Now the chapter's two weightiest voices: disagreements among \textbf{researchers themselves}, not merely insiders and outsiders.

\textbf{First: All religions are essentially the same.} Strengthen this often benevolent, evidence-bearing view. Ethical centres look similar: Jesus's instruction to treat others as one wishes to be treated in Luke and Confucius's prohibition on imposing the unwanted in ``Duke Ling of Wei,'' positive and negative forms of one structure. Practices recur beyond easy coincidence: beads, fasting in Ramadan, Buddhist afternoon abstinence, Christian fasts; pilgrimages to Mecca, Wutai, Jerusalem; confession, giving. And all ask about suffering, death, good living. Different routes to one destination seems factual rather than inflated.

\textbf{Second: Religions differ entirely; comparison itself does violence.} Its reasons are strong too. Surface likeness collapses deeper. Nirvana and Christian salvation may sound like final destinations, yet one extinguishes precisely the self's craving for endless existence while the other saves that self: nearly opposite directions. Calling both liberation narrates different games as one. Comparison is not neutral: who chooses categories and vocabulary? Europeans measured the world with Christian-derived god, Bible, church, belief, finding Asian traditions deficient, backward, superstitious. A comparison chart often disguises a report card. Finally, differences matter supremely to believers. Muslims and Christians may share bead shapes yet both object when told they believe the same thing. What authorises researchers to erase their distinctions?

Both hold truth and exceed reach. Universal unity erases what believers treasure, an exercise of power. Incomparability goes too far: beads, fasting, pilgrimage are real repetitions, and millennia of Buddhist movement from India to China and Islam into Quanzhou are historical. Mutual learning requires something comparable.

Where should we stand? First take a usable tool.

\section[Thinking Bridge: Seven Measuring Scales]{Thinking Bridge: Seven Scales, One Dimension at a Time}
Definitions failed because we demanded a black-or-white boundary: is X religion? The missing boundary breeds quarrels.

Religious-studies scholar Ninian Smart suggests, broadly, replacing a boundary with scales. Religion is a bundle of dimensions, phenomena scoring strongly or weakly on each. Combining our opening clues, we use seven:

\textbf{Belief.} Fundamental claims: gods, cosmic origins, death. Buddhist dependent arising, Christian creation, Daoist transformations of qi occupy this scale.

\textbf{Ritual.} Repeated prescribed acts: prayer, worship, incense, baptism, divination blocks, fasting. Often more central than belief; people unable to state doctrine know how to honour ancestors at New Year.

\textbf{Experience.} Awe, peace, ecstasy, forgiveness, meditative concentration. Religion is not all external rules. Many retain faith because an experience felt undeniably real.

\textbf{Community.} Do people recognise one another as insiders through this practice? Churches, sanghas, temple-fair neighbours, fasting families. Solitary convictions and communal religion are different social forms.

\textbf{Scripture.} Authoritative writing or oral tradition: Quran, Bible, Daodejing, Buddhist sutras, inherited myths.

\textbf{Ethics.} Rules for being human: Ten Commandments, Five Precepts, not imposing the unwanted, caste conventions, halal food requirements.

\textbf{Institutions.} Organisation, professionals, property, power: papacy, temple lands, imam qualifications, episcopal appointments.

The old yes-or-no problems now become readings, removing much dispute.

Confucianism? High ethics, scripture, institutions, lower experience and god-belief. Unlike Christianity, but unlike purely private philosophy too. Identify dimensions instead of forcing a verdict.

Forwarding lucky carp? Some ritual, the prescribed forwarding; vague auspicious belief; the other five nearly zero. Certainly not religion, but under its conceptual sky, addressing a related source: anxiety before uncertainty.

Astrology, yoga, qigong, fandom? Measure for yourself. Different yet overlapping answers match reality.

The real benefit is not a cleverer judgement but \textbf{freedom from needing yes-or-no judgements}. Religious study begins by giving scales the judge's seat.

\section{How Confucius Discussed Sacrifice}
Now primary material. Over two millennia ago, a teacher was directly asked our sharpest question: do spirits exist? His manner of answering remains a model.

``Advanced in Practice'' in the \textit{Analects} records:

\begin{quote}
Jilu asked about serving spirits. The Master said: ``Unable yet to serve people, how can you serve spirits?'' He ventured to ask about death. ``Not yet knowing life, how can you know death?''
\end{quote}
Broadly, Zilu asks how to serve spirits; Confucius redirects to serving people. Asked about death, he redirects to understanding life.

``Eight Rows'' adds:

\begin{quote}
Sacrifice as though present; sacrifice to spirits as though they were present. The Master said: ``If I do not participate, it is as though I had not sacrificed.''
\end{quote}
Honour ancestors as really present, spirits likewise. Personal nonparticipation, even with a substitute, amounts to no rite.

Read together and see the precision. Zilu supposedly asks existence questions, expecting yes or no about spirits and death. Confucius chooses neither, turning entirely: \textbf{before asking whether spirits are there, consider who you are before the ritual table}. Sacrifice as though present is subtle: not because deity exists, but as though it does. Attention moves from deity to worshipper's sincerity, concentration, personal presence, with delegation insufficient.

Notice what he \textbf{does not do}. He neither mocks superstition, already an available Spring and Autumn attitude, nor invents myths proving heavenly ancestors, plentiful elsewhere. He surgically separates ritual meaning from supernatural truth, addressing the former and suspending the latter.

This distinction leads the chapter's second half. Today's scholars use the same structure: fully describe a rite's participant meaning without deciding divine existence. Conversely, even believing gods nonexistent cannot make a worshipper's sincere reverence unreal. Both outside and inside observers can recognise that actual human experience.

\section{One Prayer, Three Explanations}
Return to the grandmother raising incense, bowing, throwing blocks. These observable events are undisputed inside or out; why opens disagreement. Compare genuinely different, nonequivalent explanations, each with reasons and limits. Keep four drawers separate: events, explanations, participant understanding, evaluation.

\textbf{First: She responds to divine guidance (the participant's account).}

Guan Yu is effective; blocks communicate his answer; the favourable throw means acknowledgement and agreement. This is indispensable first-hand testimony about understanding, not automatically the best scientific cause. Skip it for she is merely anxious and you study an invented woman. Its limit: no publicly shared outside test settles efficacy between believers and sceptics. Its strength: it comprehensively explains her next acts, calmer feelings, soup for her grandson, next year's return.

\textbf{Second: Ritual manages uncontrollable anxiety (psychological explanation).}

Early twentieth-century anthropologist Bronislaw Malinowski noted a famous Trobriand Islands contrast. Calm-lagoon fishing involved little magic; open-sea fishing amid storms, currents, mortal uncertainty required meticulous ritual. Broadly, magic and religion arise where technique ends and risk begins, supporting the feeling of having done everything possible. Grandma has cooked, advised, arranged tutoring; luck now controls what cannot be made into soup. Blocks address that remainder. This cross-cultural account also locates your exam-eve carp. Its limit: why Guan Yu rather than another ritual, and what about gratitude after fulfilled vows or respect for loyalty, unrelated to anxiety? Calling everything anxiety relief resembles reducing a love song to vibrating vocal cords: not wrong, but missing the song.

\textbf{Third: Ritual binds community (sociological explanation).}

In \textit{The Elementary Forms of Religious Life} (1912), Emile Durkheim proposed, broadly, that worshipping society honours itself; recurring ritual joins separate individuals into us. This temple is not merely a private relationship with Guan Yu. It is fairs, neighbouring joss-paper shops, communal calendars, returning emigrants' destination. Guan Yu's righteousness expresses merchant society's needed loyalty and trust. A trading city annually reaffirms its foundations. This explains ritual publicity and temples' local centrality beyond the first accounts. Its limit: translating everything into social need neglects solitary midnight incense; silence about Guan Yu's actual existence leaves believers feeling the question unanswered.

These explanations occupy different floors: participant world, individual mind, collective structure. They clash chiefly when each claims completeness. Mature researchers identify which door each key opens, not that only theirs counts as a key.

\section[Further Challenge: Suspend Judgement]{Further Challenge: Suspend Judgement, Not Your Mind}
Now the central methodological issue, professionally named but already demonstrated by Confucius: \textbf{temporarily bracket the truth of supernatural claims when studying religion}.

Bracketing is not deletion. Public-evidence research uses texts, archaeology, interviews, statistics, comparison to investigate beliefs, actions, social effects, not establish God's existence. Responsible phrasing distinguishes a tradition claiming resurrection, believers trusting divine replies through blocks, and historical evidence dating a mosque to Northern Song. Three verbs mark three levels. Confusing them misrepresents: historical evidence proves God and believers imagined the whole tradition equally overreach, in opposite directions.

Criticism comes from both sides and deserves serious replies.

Insiders ask how outsiders can understand. Which is truer, participants' experiences, language, reasons, or observers' descriptions and comparisons? Without meditation or conversion's shattering experience, writing merely scratches through a shoe. This has weight: experience is difficult to acquire from paper. Yet proximity also obscures. Inside Mount Lu one cannot see its full face. Grandma may not recognise her rite's structure in Trobriand fishing, as fish do not notice water. Comparison, distance, outside vocabulary illuminate uniquely. We need both legs. Good description lets Grandma recognise herself and an uninitiated reader understand why she acts, translating between perspectives rather than privileging one alone.

The other criticism calls bracketing sophisticated evasion. If remedies kill a child or a leader enriches himself through divine authority, does suspended judgement make scholarship complicit? Accept the challenge fully: it identifies the method's real boundary. Bracket \textbf{supernatural truth}, not \textbf{human consequences}. God's existence is outside the tools' reach; whose pocket receives money, whether a child receives medical care, whose freedom is removed are publicly investigable. Researchers share every citizen's responsibility here. Understanding meaning never exempts deeds under its banner; explanation never equals defence, as the previous chapter's tool taught.

Another counterexample prevents lazy bracketing. A supposedly devout believer clearly declares affiliation, knows doctrine, attends organisation, rejects rivals. Apply that template to Japan, where religious rites densely stitch lives: shrine blessings for infants, church or shrine weddings, Buddhist year-end bells, almost universal monastic funerals, yet surveys often find fewer than thirty percent declaring religion. The template calls them irreligious and rites empty. Tears at funerals and joined hands at shrines resist that verdict. The template grew from confessed doctrine and church membership. Applied to religion living in custom rather than declaration, the blank belongs to the ruler, not society. Bracket not only gods' truth but your own scale; inspect before measuring.

Gather this through Wittgenstein's family resemblance. Games include boards, cards, balls, hide-and-seek without one feature shared by all and only games. Like relatives' eyes and chins, likeness links without one trait spanning everyone. Religion works similarly: Buddhist doctrine with less deity, Shinto ritual with fewer scriptures, Confucian ethics with less experiential emphasis. Overlap weaves a net without one end-to-end thread. Hence neither sameness nor incomparability holds, and a clean boundary pursues fish in a tree. Take seven scales and meet each family member individually; that is sufficient.

\section{The Age of Forwarded Lucky Carp}
The old question wears countless phone disguises daily.

Each June exam season peaks in cyber-incense: forward carp or Manjushri, comment with pleas to pass. Electronic wooden-fish apps reach millions of downloads: tap once, gain merit. Astrologers out-follow many subject teachers; tarot streams run all night. Opposing voices are equally loud: how can you believe this now, astrology is a gullibility tax, practice papers beat lucky forwarding.

Our tools show both sides mostly punching air.

Measure carp. Belief is weak and vague; few seriously posit a carp deity, more a whisper to luck. Ritual is standard: forwarding with blessing, before results, prohibitions such as liking one's own post. Experience offers cheap reassurance; the other four scales approach zero. Not religion, a point both can accept. Yet its function belongs beside blocks and fishing rites, supplying agency where effort ends and luck begins. Mockers may themselves avoid fourth floors---four sounds like death in Chinese---or certain exam clothes. Seeing this protects not carp but perceptiveness about people.

Astrology adds dimensions: systematic chart literature, community through zodiac identities, ethical guidance such as this week's decluttering, even the shiver of feeling accurately described. Institutions and serious supernatural commitment are weaker. Seven scales locate its distance and overlap with traditional religion more usefully than an insult, also asking why intelligent people need something unscientific. Revisit anxiety, then consider the rarity of feeling precisely understood.

Fandom called a cult can likewise be analysed: strong community, elaborate chart-boosting, support and comment-control rituals, explicit loyalty ethics, some canon and institution, zero supernatural belief. Readings upgrade abuse into analysis: where are similarities, which youthful needs do intense dimensions meet, and what else might house those needs?

Suspending judgement is daily work. A headscarf-wearing deskmate invites three distinctions: understand why through listening and reading; recognise a seriously chosen life through respect; judge and speak if anyone harms her under faith's name, without respect silencing criticism. Reformation, Crusades, Buddhist suppression and patronage in history similarly belong in four drawers: events, explanations, participant meanings, evaluations. Neither arrogant judge nor bland bystander is necessary.

Electronic wooden fish and AI divination push further. If a program replaces even a human at ritual's other end, what remains of a merit-adding tap? Perhaps nothing, perhaps ritual's core. Confucius already locates the crucial point with the worshipper, not the worshipped: sacrifice as though present. Whether that supports a digital fish is yours to weigh.

\section{Your Turn: How Wide Can the Brackets Be?}
Return to the Quanzhou street and Xiaolin's face.

He now knows what superstition sweeps away, how seven scales work, the precision of as though present, and that bracketing supernatural claims does not bracket human good and evil. More tools make one question heavier, without a standard answer.

Where should methodological bracketing end?

One side is clear. Without it, research becomes sentencing. Declare Guan Yu nonexistent and Grandma's door closes, losing not only courtesy but understanding of substantial human life. History repeatedly shows judges certain of ultimate truth smashing more than clay images.

The other is clear too. Oversized brackets become indulgence. A master bankrupts devotees, doctrine forbids girls' schooling, a group condemns blasphemers to death. Declaring oneself only a researcher becomes an ignoble position. Someone must reckon human accounts.

Where draw understanding versus judgement? Believe anything without violence? Intervene once money is deceived away? Does manipulative fear-marketing count? Different rules for minors? What about adults' willing offerings? Should scholars, law, or each informed ordinary person repeatedly draw the boundary?

Confucius suspended spirits' existence yet judged human right and wrong throughout life. These did not conflict. Perhaps an answer lies in their distance; perhaps not.

Where would you draw the line? Give reasons, then ask whether those reasons themselves can be bracketed and examined.
